\definecolor{taxblue}{RGB}{38, 87, 139}
\definecolor{taxlight}{RGB}{238, 244, 251}
\definecolor{taxline}{RGB}{70, 70, 70}

\resizebox{0.96\linewidth}{!}{%
\begin{tikzpicture}[
  font=\sffamily,
  >=Latex,
  layer/.style={
    draw=taxblue,
    rounded corners=4pt,
    line width=0.7pt,
    fill=taxlight,
    align=center,
    minimum height=0.86cm,
    text width=3.25cm,
    font=\sffamily\small\bfseries,
  },
  note/.style={
    draw=taxblue,
    rounded corners=3pt,
    line width=0.55pt,
    fill=white,
    align=center,
    text width=3.25cm,
    minimum height=0.72cm,
    font=\sffamily\scriptsize,
  },
  arrow/.style={->, line width=0.55pt, draw=taxline},
]

\node[layer] (core) at (0,0) {Core failure label};
\node[note, below=0.18cm of core] (corenote) {dominant reason not to show as-is};

\node[layer, right=0.62cm of core] (modifier) {Secondary modifier};
\node[note, below=0.18cm of modifier] (modifiernote) {additional detail for interpretation or mitigation};

\node[layer, right=0.62cm of modifier] (dimension) {Evaluation dimensions};
\node[note, below=0.18cm of dimension] (dimensionnote) {correctness, grounding, usefulness, actionability};

\node[layer, right=0.62cm of dimension] (decision) {Mitigation decision};
\node[note, below=0.18cm of decision] (decisionnote) {show, rewrite, suppress, or escalate};

\node[layer, right=0.62cm of decision] (confidence) {Confidence label};
\node[note, below=0.18cm of confidence] (confidencenote) {annotator certainty under available context};

\draw[arrow] (core) -- (modifier);
\draw[arrow] (modifier) -- (dimension);
\draw[arrow] (dimension) -- (decision);
\draw[arrow] (decision) -- (confidence);

\node[below=0.55cm of dimensionnote, font=\sffamily\footnotesize\itshape, text=taxline]
  {Failure labels explain the problem; decision labels define the action used in trade-off analysis.};

\end{tikzpicture}%
}
